\documentclass[11pt,a4paper]{article}

% ── Packages ──────────────────────────────────────────────────────────────────
\usepackage[margin=1in]{geometry}
\usepackage{amsmath}
\usepackage{booktabs}
\usepackage{graphicx}
\usepackage{hyperref}
\usepackage{caption}
\usepackage{subcaption}
\usepackage{listings}
\usepackage{xcolor}
\usepackage{microtype}
\usepackage{parskip}
\usepackage{enumitem}
\usepackage{siunitx}
\usepackage{fancyhdr}
\usepackage{titling}

% ── Hyperref setup ────────────────────────────────────────────────────────────
\hypersetup{
    colorlinks=true,
    linkcolor=blue!60!black,
    citecolor=green!50!black,
    urlcolor=blue!60!black
}

% ── Code listing style ────────────────────────────────────────────────────────
\lstset{
    basicstyle=\ttfamily\small,
    keywordstyle=\color{blue!70!black}\bfseries,
    commentstyle=\color{gray},
    stringstyle=\color{orange!80!black},
    backgroundcolor=\color{gray!8},
    frame=single,
    rulecolor=\color{gray!40},
    breaklines=true,
    captionpos=b,
    numbers=left,
    numberstyle=\tiny\color{gray},
    xleftmargin=1.5em,
    framexleftmargin=1.5em,
    language=bash
}

% ── Header / Footer ───────────────────────────────────────────────────────────
\pagestyle{fancy}
\fancyhf{}
\rhead{CA Assignment — TSP \& Knapsack}
\lhead{Abimanyu \& Yuvraj}
\cfoot{\thepage}

% ─────────────────────────────────────────────────────────────────────────────
\begin{document}

% ── Title Page ────────────────────────────────────────────────────────────────
\begin{titlepage}
    \centering
    \vspace*{2cm}
    {\LARGE\bfseries Performance Analysis of Dynamic Programming\\[0.4em]
    Using Linux perf and ChampSim\par}
    \vspace{1.5cm}
    {\large Abimanyu Choudhary K (106124006) \quad Yuvraj Meena (106124142)\par}
    \vspace{0.8cm}
    {\large \today\par}
    \vspace{2cm}
    \begin{abstract}
    \noindent
    In this project we evaluated CPU and memory performance of two combinatorial
    optimisation algorithms implemented as dynamic-programming workloads in C.
    We used two complementary tools: \emph{Linux perf} for native hardware-counter
    profiling on our physical processor, and \emph{ChampSim} for trace-driven
    microarchitectural simulation.  The two workloads are the \emph{Traveling
    Salesman Problem} solved by the Held--Karp bitmask-DP algorithm
    ($\mathcal{O}(2^N N^2)$ time, $\mathcal{O}(2^N N)$ space), and the \emph{0/1
    Knapsack Problem} solved by a bottom-up 1-D rolling-array DP
    ($\mathcal{O}(nW)$ time).  TSP exercises irregular, random-access memory with
    a ${\sim}80$\,MiB working set that spills past all cache levels, while
    Knapsack uses a backward-sweep over a~${\sim}1.9$\,MiB array that fits in
    L2/L3 but not L1D.  We measured instructions, cycles, IPC, cache-miss rates,
    and branch-prediction accuracy.  For the simulation component, we built two
    ChampSim binaries with L1D cache capacities of 32\,KiB (32 sets $\times$ 8-way)
    and 64\,KiB (64 sets $\times$ 8-way), and replayed the standard
    \texttt{400.perlbench-41B} instruction trace.  This report details our
    methodology, experimental results, analysis, and full reproduction steps.
    \end{abstract}
\end{titlepage}

\tableofcontents
\newpage

% ─────────────────────────────────────────────────────────────────────────────
\section{Introduction and Objectives}
% ─────────────────────────────────────────────────────────────────────────────

When optimising software or designing processor hardware, relying on wall-clock
time alone rarely tells the full story.  A program can run slowly because of
back-end stalls from cache misses, front-end stalls from branch mispredictions,
or memory-bandwidth saturation.

In this lab our goals were:

\begin{itemize}[leftmargin=*]
    \item \textbf{Linux perf}: Inspect real hardware counters on our machine and
          observe how contrasting DP algorithms with different memory-access
          patterns affect IPC, cache references, and branch behaviour.
    \item \textbf{ChampSim}: Isolate microarchitectural components in a
          simulator and observe how changing L1 data-cache capacity affects
          program execution under controlled trace replays.
    \item \textbf{Reproducibility}: Document configurations, source code, exact
          shell commands, and terminal outputs so every experiment can be
          independently verified.
\end{itemize}

% ─────────────────────────────────────────────────────────────────────────────
\section{Background}
% ─────────────────────────────────────────────────────────────────────────────

\subsection{Measuring Hardware Events with Linux perf}

Modern processors contain dedicated hardware units called \emph{Performance
Monitoring Units} (PMUs).  The PMU maintains special hardware registers that
increment on microarchitectural events such as retired instructions, elapsed
core clock cycles, cache hits and misses, and mispredicted branches~\cite{gregg2020}.

The primary metric we use to gauge execution throughput is
\emph{Instructions Per Cycle} (IPC):

\begin{equation}
    \text{IPC} = \frac{\text{Instructions Retired}}{\text{Cycles Elapsed}}
    \label{eq:ipc}
\end{equation}

A higher IPC indicates that the CPU's out-of-order execution engine is
successfully exploiting instruction-level parallelism (ILP).  Conversely,
frequent L1 data-cache misses or branch mispredictions stall the pipeline and
drag IPC down~\cite{hp2017}.

\subsection{Trace-Driven Simulation with ChampSim}

While \texttt{perf} measures physical execution on a real CPU (with all its OS
noise and prefetchers), ChampSim simulates an idealised cycle-level CPU model
by replaying an instruction trace~\cite{hp2017}.  This allows us to modify
architectural parameters—such as cache sets, associativity, or replacement
policies—that are impossible to change on physical silicon.

% ─────────────────────────────────────────────────────────────────────────────
\section{Experimental Setup}
% ─────────────────────────────────────────────────────────────────────────────

\subsection{Test System Specifications}

All native benchmarks were run on a Linux machine with the following
specifications:

\begin{itemize}[leftmargin=*]
    \item \textbf{Processor}: \textit{[fill in from your VM / machine]}
    \item \textbf{Cores / Threads}: \textit{[fill in]}
    \item \textbf{Clock Frequency}: \textit{[fill in]}
    \item \textbf{Compiler}: GCC with optimisation flag \texttt{-O2}
    \item \textbf{Simulation}: ChampSim built from source with custom JSON
          cache configurations
\end{itemize}

\subsection{Selected Native Workloads}

We chose two NP-hard combinatorial optimisation problems that are canonically
solved with dynamic programming, but exhibit fundamentally different memory
traversal patterns:

\begin{enumerate}[leftmargin=*]

    \item \textbf{Traveling Salesman Problem — Held--Karp DP}
          ($N = 20$ cities).
          This algorithm maintains a \emph{bitmask DP table}
          $\mathrm{dp}[S][i]$, where $S \subseteq \{0,\ldots,N-1\}$ is a
          subset of visited cities (represented as a 20-bit integer) and $i$
          is the last city visited.  The full table has $2^{20} \times 20$
          entries (${\approx}80$\,MiB for two \texttt{int} arrays).
          Because the transition
          $\mathrm{dp}[\,S\cup\{v\}\,][v] \leftarrow \mathrm{dp}[S][u] +
          \mathrm{dist}[u][v]$
          jumps between arbitrary bitmask states, the access pattern is
          highly \emph{irregular}—stressing the TLB and producing frequent
          L1/L2 capacity misses, similar to MCM from the reference report
          but at larger scale.

    \item \textbf{0/1 Knapsack Problem — rolling-array DP}
          ($n = 2{,}000$ items, $W = 500{,}000$).
          We use the standard single-pass \emph{backward} sweep:
          $\mathrm{dp}[w] \leftarrow \max(\mathrm{dp}[w],\;
          \mathrm{dp}[w - w_i] + v_i)$ for $w = W$ down to $w_i$.
          The \texttt{dp[]} array occupies
          $(W+1) \times 4 \approx 1.9$\,MiB—comfortably inside a modern
          L2/L3 cache but exceeding L1D, so the hardware prefetcher can
          bring lines ahead of the backward sweep, producing intermediate
          cache behaviour between the two reference workloads.
\end{enumerate}

\subsection{ChampSim Simulation Configuration}

For the simulator component we replayed the standard SPEC CPU benchmark trace
\texttt{400.perlbench-41B.champsimtrace.xz}.  Following standard simulation
practices, we configured a warm-up period of 100{,}000 instructions to
populate the caches and branch predictor, followed by 1{,}000{,}000 detailed
simulation instructions.

We designed two configurations where \emph{only} the L1D cache capacity
differed:

\begin{itemize}[leftmargin=*]
    \item \textbf{32\,KiB Control}: 32 sets $\times$ 8-way $\times$ 64\,B
          lines $= 16{,}384$ bytes, i.e.\ 32\,KiB.
    \item \textbf{64\,KiB Experiment}: 64 sets $\times$ 8-way $\times$ 64\,B
          lines $= 32{,}768$ bytes, i.e.\ 64\,KiB.
\end{itemize}

Both builds shared identical replacement policies (LRU), block sizes (64\,B),
and associativity (8-way), so any observed difference is attributable solely
to cache capacity.

% ─────────────────────────────────────────────────────────────────────────────
\section{Native Hardware Results (Linux perf)}
% ─────────────────────────────────────────────────────────────────────────────

\subsection{TSP — Held--Karp Analysis}

We executed the compiled TSP binary under \texttt{perf stat} with 3
repetitions.  The recorded hardware-counter metrics are listed in
Table~\ref{tab:tsp}.

\begin{table}[h!]
\centering
\caption{Hardware performance counters for TSP (Held--Karp, $N=20$)}
\label{tab:tsp}
\begin{tabular}{lS[table-format=3.3e3]}
\toprule
\textbf{Counter / Metric} & {\textbf{Recorded Value}} \\
\midrule
Instructions retired         & {[FILL FROM tsp\_perf.txt]} \\
Cycles elapsed               & {[FILL FROM tsp\_perf.txt]} \\
Instructions Per Cycle (IPC) & {[FILL FROM tsp\_perf.txt]} \\
Cache references             & {[FILL FROM tsp\_perf.txt]} \\
Cache misses                 & {[FILL FROM tsp\_perf.txt]} \\
Cache miss rate              & {[FILL \%]} \\
Total branches               & {[FILL FROM tsp\_perf.txt]} \\
Branch misses                & {[FILL FROM tsp\_perf.txt]} \\
Branch miss rate             & {[FILL \%]} \\
Average run time             & {[FILL seconds]} \\
Timing variance              & {[FILL \%]} \\
\bottomrule
\end{tabular}
\end{table}

Key observations from the TSP data:

\begin{itemize}[leftmargin=*]
    \item \textbf{IPC}: The irregular bitmask-indexed access pattern prevents
          the hardware prefetcher from loading data ahead of demand.
          Every state transition
          $\mathrm{dp}[S][u] \to \mathrm{dp}[S \cup \{v\}][v]$ reads from a
          memory offset proportional to the bitmask $S$, which sweeps in
          Gray-code-like order—not sequential.  This results in persistent
          L2/L3 misses that throttle the pipeline.

    \item \textbf{Cache miss rate}: The 80\,MiB working set of the TSP DP
          table far exceeds both L1D and L2 caches on most modern processors.
          We expect a significantly higher miss rate than Knapsack.

    \item \textbf{Branch miss rate}: The inner loop has a predictable loop
          bound (iterating over bits 0 to $N-1$), so the branch predictor
          handles the fixed-count loop well despite the irregular data access.

    \item \textbf{Page faults}: Allocating the ${\sim}80$\,MiB DP table via
          \texttt{malloc} causes Linux to fault in memory pages on first
          access, contributing to startup overhead.
\end{itemize}

\subsection{Knapsack — 0/1 DP Analysis}

Next we evaluated the Knapsack benchmark under the same 3-run configuration.
Results are summarised in Table~\ref{tab:ks}.

\begin{table}[h!]
\centering
\caption{Hardware performance counters for 0/1 Knapsack ($n=2000$, $W=500000$)}
\label{tab:ks}
\begin{tabular}{lS[table-format=3.3e3]}
\toprule
\textbf{Counter / Metric} & {\textbf{Recorded Value}} \\
\midrule
Instructions retired         & {[FILL FROM knapsack\_perf.txt]} \\
Cycles elapsed               & {[FILL FROM knapsack\_perf.txt]} \\
Instructions Per Cycle (IPC) & {[FILL FROM knapsack\_perf.txt]} \\
Cache references             & {[FILL FROM knapsack\_perf.txt]} \\
Cache misses                 & {[FILL FROM knapsack\_perf.txt]} \\
Cache miss rate              & {[FILL \%]} \\
Total branches               & {[FILL FROM knapsack\_perf.txt]} \\
Branch misses                & {[FILL FROM knapsack\_perf.txt]} \\
Branch miss rate             & {[FILL \%]} \\
Page faults                  & {[FILL FROM knapsack\_perf.txt]} \\
Average run time             & {[FILL seconds]} \\
Timing variance              & {[FILL \%]} \\
\bottomrule
\end{tabular}
\end{table}

Comparing Knapsack to TSP reveals clear microarchitectural differences:

\begin{itemize}[leftmargin=*]
    \item \textbf{IPC}: The 1-D \texttt{dp[]} array of ${\sim}1.9$\,MiB fits
          entirely in L2/L3.  Although the backward sweep is not as prefetcher-
          friendly as a pure forward scan, the stride is fixed and the hardware
          prefetcher can still anticipate accesses for small item weights,
          keeping the pipeline better fed than in TSP.

    \item \textbf{Cache miss rate}: Because the working set fits in L3, the
          sustained miss rate to main memory is much lower than for TSP,
          despite the backward access direction.

    \item \textbf{Near-zero branch mispredictions}: The two nested \texttt{for}
          loops have fixed bounds, making the branch pattern trivially
          predictable.  The only misprediction source is the conditional
          \texttt{if (candidate > dp[w])}.

    \item \textbf{Instruction count}: Because Knapsack iterates $n \times W =
          2000 \times 500{,}000 = 10^9$ inner steps, the total retired
          instruction count will be substantially higher than TSP, which
          iterates $2^{20} \times 20^2 \approx 4 \times 10^8$ transitions.
\end{itemize}

\subsection{Comparative Summary (native perf)}

\begin{table}[h!]
\centering
\caption{Side-by-side comparison of TSP and Knapsack hardware counters}
\label{tab:compare}
\begin{tabular}{lcc}
\toprule
\textbf{Metric} & \textbf{TSP (Held--Karp)} & \textbf{Knapsack (0/1 DP)} \\
\midrule
IPC                  & ---   & ---   \\
Cache miss rate (\%) & ---   & ---   \\
Branch miss rate (\%)& ---   & ---   \\
Run time (s)         & ---   & ---   \\
\bottomrule
\end{tabular}
\end{table}

% Insert figures after running plot_results.py
\begin{figure}[h!]
    \centering
    \includegraphics[width=0.55\textwidth]{../results/fig_ipc_bar.pdf}
    \caption{Native IPC comparison between TSP and Knapsack workloads.}
    \label{fig:ipc}
\end{figure}

\begin{figure}[h!]
    \centering
    \begin{subfigure}[b]{0.48\textwidth}
        \includegraphics[width=\textwidth]{../results/fig_cache_bar.pdf}
        \caption{Cache miss rate comparison.}
        \label{fig:cache}
    \end{subfigure}
    \hfill
    \begin{subfigure}[b]{0.48\textwidth}
        \includegraphics[width=\textwidth]{../results/fig_branch_bar.pdf}
        \caption{Branch miss rate comparison.}
        \label{fig:branch}
    \end{subfigure}
    \caption{Cache and branch predictor miss rates for TSP vs.\ Knapsack.}
    \label{fig:miss_rates}
\end{figure}

% ─────────────────────────────────────────────────────────────────────────────
\section{ChampSim Simulation Results}
% ─────────────────────────────────────────────────────────────────────────────

We ran our two custom ChampSim binaries on the Perlbench trace and logged
instructions, cycles, IPC, and cache misses.  The comparative metrics are
shown in Table~\ref{tab:champsim}.

\begin{table}[h!]
\centering
\caption{ChampSim simulation results — 32\,KiB vs.\ 64\,KiB L1D}
\label{tab:champsim}
\begin{tabular}{lcccc}
\toprule
\textbf{Configuration} & \textbf{Instructions} & \textbf{Cycles} &
\textbf{IPC} & \textbf{L1D Misses} \\
\midrule
32\,KiB Control (32 sets) & --- & --- & --- & --- \\
64\,KiB Experiment (64 sets) & --- & --- & --- & --- \\
\bottomrule
\end{tabular}
\end{table}

\begin{figure}[h!]
    \centering
    \begin{subfigure}[b]{0.48\textwidth}
        \includegraphics[width=\textwidth]{../results/fig_champsim_ipc.pdf}
        \caption{IPC comparison: 32\,KiB vs.\ 64\,KiB L1D.}
        \label{fig:cs_ipc}
    \end{subfigure}
    \hfill
    \begin{subfigure}[b]{0.48\textwidth}
        \includegraphics[width=\textwidth]{../results/fig_champsim_misses.pdf}
        \caption{L1D miss count comparison.}
        \label{fig:cs_misses}
    \end{subfigure}
    \caption{ChampSim simulation results for the Perlbench trace.}
    \label{fig:champsim}
\end{figure}

\subsection{Analysis of Simulation Results}

\textit{[Fill in after running ChampSim.  Follow the pattern from the reference
report: if both configurations give identical results, explain working-set
size.  If they differ, explain the capacity-miss reduction.]}

The standard explanation for identical miss counts across two cache sizes is
architectural: if the program's active working set fits entirely within the
smaller configuration, increasing capacity adds empty sets that are never
used.  Increasing cache size only reduces misses when the working set
\emph{exceeds} the current cache capacity~\cite{hp2017}.

% ─────────────────────────────────────────────────────────────────────────────
\section{Discussion and Lessons Learned}
% ─────────────────────────────────────────────────────────────────────────────

\begin{itemize}[leftmargin=*]
    \item \textbf{Memory-access pattern is the dominant bottleneck.}
          TSP's random bitmask indexing and large working set expose the
          processor to sustained L2/L3 (and potentially DRAM) latency,
          far outweighing the algorithmic instruction-count advantage it
          might have over Knapsack.

    \item \textbf{Working-set size vs.\ cache hierarchy.}
          Knapsack's ${\sim}1.9$\,MiB rolling array fits comfortably in L2
          or L3 on modern processors.  Even though the backward sweep is
          slightly less prefetcher-friendly than a forward scan, the cache
          hit rate is dramatically higher than TSP.

    \item \textbf{Native profiling vs.\ simulation.}
          Linux \texttt{perf} gives ground-truth numbers on live silicon but
          includes non-deterministic OS noise.  ChampSim gives perfectly
          repeatable, deterministic cycle counts, but only reflects the
          specific trace snippet being simulated.

    \item \textbf{Branch predictability.}
          Both algorithms use nested loops with fixed bounds, making the
          branch predictor's job straightforward.  NP-hard algorithmic
          complexity does not imply branch unpredictability.
\end{itemize}

% ─────────────────────────────────────────────────────────────────────────────
\section{Conclusion}
% ─────────────────────────────────────────────────────────────────────────────

Through this project we successfully collected and analysed hardware-counter
statistics for two classic NP-hard optimisation problems solved with dynamic
programming, using Linux \texttt{perf} for native measurements and ChampSim
for trace-driven simulation.

Our measurements highlight that the \emph{memory-access pattern} of a DP
algorithm—not its time complexity—is the dominant factor governing CPU
performance.  TSP's irregular bitmask indexing over an 80\,MiB table leads
to persistent cache pressure, while Knapsack's backward sweep over a
1.9\,MiB array enables much more effective cache reuse.  The ChampSim
experiment illustrates the classic cache-design principle: increasing cache
capacity is only effective when the program's working set exceeds the current
boundary.

% ─────────────────────────────────────────────────────────────────────────────
\appendix
\section{Execution Manual}
% ─────────────────────────────────────────────────────────────────────────────

This appendix records the exact commands used to run and reproduce all
experiments.  See \texttt{SETUP.md} in the project root for the complete
environment setup guide.

\subsection{Virtual Environment Setup}
\begin{lstlisting}
cd ~/perf-champsim-project
python3 -m venv .venv
source .venv/bin/activate
pip install matplotlib pandas
\end{lstlisting}

\subsection{Compiling the C Workloads}
\begin{lstlisting}
cd ~/perf-champsim-project
gcc -O2 perf/tsp.c      -o bin/tsp
gcc -O2 perf/knapsack.c -o bin/knapsack
\end{lstlisting}

\subsection{Profiling TSP with perf}
\begin{lstlisting}
cd ~/perf-champsim-project
perf stat -r 3 \
  -e task-clock,context-switches,cpu-migrations,page-faults,\
cycles,instructions,branches,branch-misses,cache-references,cache-misses \
  ./bin/tsp 2>&1 | tee results/tsp_perf.txt
\end{lstlisting}

\subsection{Profiling Knapsack with perf}
\begin{lstlisting}
cd ~/perf-champsim-project
perf stat -r 3 \
  -e task-clock,context-switches,cpu-migrations,page-faults,\
cycles,instructions,branches,branch-misses,cache-references,cache-misses \
  ./bin/knapsack 2>&1 | tee results/knapsack_perf.txt
\end{lstlisting}

\subsection{Running ChampSim — 32\,KiB Control}
\begin{lstlisting}
cd ~/perf-champsim-project
./bin/champsim_l1d_32kb \
  --warmup-instructions 100000 \
  --simulation-instructions 1000000 \
  --json results/perlbench_32kb.json \
  traces/400.perlbench-41B.champsimtrace.xz \
  2>&1 | tee results/perlbench_32kb.txt
\end{lstlisting}

\subsection{Running ChampSim — 64\,KiB Experiment}
\begin{lstlisting}
cd ~/perf-champsim-project
./bin/champsim_l1d_64kb \
  --warmup-instructions 100000 \
  --simulation-instructions 1000000 \
  --json results/perlbench_64kb.json \
  traces/400.perlbench-41B.champsimtrace.xz \
  2>&1 | tee results/perlbench_64kb.txt
\end{lstlisting}

\subsection{Generating Plots}
\begin{lstlisting}
cd ~/perf-champsim-project
source .venv/bin/activate
python3 plot_results.py
\end{lstlisting}

% ─────────────────────────────────────────────────────────────────────────────
\section{Terminal Screenshot Evidence}
% ─────────────────────────────────────────────────────────────────────────────

\textit{Insert terminal screenshots here after running experiments.  Capture
screenshots of:}
\begin{itemize}
    \item TSP \texttt{perf stat} terminal output
    \item Knapsack \texttt{perf stat} terminal output
    \item ChampSim execution output for 32\,KiB and 64\,KiB runs
    \item Cache-geometry verification (\texttt{champsim\_config\_32kb.json}
          vs.\ \texttt{champsim\_config\_64kb.json})
\end{itemize}

% ─────────────────────────────────────────────────────────────────────────────
\begin{thebibliography}{9}
\bibitem{gregg2020}
B.\ Gregg, \textit{Systems Performance: Enterprise and the Cloud}, 2nd ed.\
Boston, MA: Addison-Wesley, 2020.

\bibitem{hp2017}
J.\ L.\ Hennessy and D.\ A.\ Patterson,
\textit{Computer Architecture: A Quantitative Approach}, 6th ed.\
Cambridge, MA: Morgan Kaufmann, 2017.

\bibitem{cormen2022}
T.\ H.\ Cormen, C.\ E.\ Leiserson, R.\ L.\ Rivest, and C.\ Stein,
\textit{Introduction to Algorithms}, 4th ed.\
Cambridge, MA: The MIT Press, 2022.

\bibitem{held1962}
M.\ Held and R.\ M.\ Karp,
``A dynamic programming approach to sequencing problems,''
\textit{Journal of the Society for Industrial and Applied Mathematics},
vol.\ 10, no.\ 1, pp.\ 196--210, 1962.

\bibitem{champsim}
ChampSim Development Team, ``ChampSim: A Trace-Based Simulator,''
GitHub repository, \url{https://github.com/ChampSim/ChampSim}, 2023.
\end{thebibliography}

\end{document}
